\section{Data, Provenance and Experimental Design}
\label{sec:data-design}
\chapterstatus{Editorial synthesis of the accepted provenance, notation and gate contracts. Provider-specific values remain governed by their source freezes.}

\subsection{Data universe and dates}

The legacy corpus contains different evidence domains: long-run gold and silver
time series, London Stock Exchange market and option inputs, mine-level
operating data and synthetic DCF examples. Each retains its own provider,
frequency, date range, unit and measure. The operating domain is quarterly and
mine-level: Team~4 estimates on 2019Q1--2025Q4 Barrick reports, with shorter
effective samples where leading observations are missing; 2026Q1--Q2 enter
as actuals and forecasts cover 2026Q3--2030Q4
(Section~\ref{sec:company-perimeter}). The 2 April 2026 GLD option snapshot
and standalone Team~4 GBM paths are a separate legacy domain, excluded from
the integrated valuation \cite{team4}. Gold histories, option-calibration
surfaces and operating panels are not pooled.

The copper extension is another declared experiment: historical HG/HXE
bid/ask surfaces, contemporaneous matching futures and date-specific Treasury
curves are distinct from both GLD and the mine panel. Its preliminary temporal
pilot samples ten dates across 9 July--2 September 2026, including the final
three consecutive dates, with approximately 100 spaced real locations per
date. It is not a reproduction of the 31-date gold OOS panel. Missing, stale
or wide-spread quotes remain exclusions; no synthetic quotes complete a date.
The surviving contract catalogue prevents a complete point-in-time universe
and expired-option history is unavailable from this API
\cite{ib-historical-limitations}. Admitted dates and actual forecast gaps are
reported in Section~\ref{sec:copper-extension}.

\subsection{Provenance and reproducibility}

Every quantitative result admitted to the final study must be traceable through an input hash, as-of date, configuration, seed where applicable, executable stage, output path and tolerance. Exact reuse identifies material preserved from a source article; edited reuse identifies rearrangement or notation harmonisation; recomputed identifies a result created by a frozen executable run. A gap never authorises a fabricated value.

\subsection{Common notation, units and model interfaces}

The notation distinguishes the physical measure $\mathbb P$ from the pricing
measure $\mathbb Q$, GLD in USD/share from gold in USD/oz, annualised decimals
from percentage points, and enterprise from equity value. Production is
$\mathrm{Prod}$ or italic $Q$ in Chapter~11, distinct from $\mathbb Q$.
Ore is in tonnes, grade in g/t, recovery decimal, and production in troy
ounces of $31.1035$ grams, usually reported in thousands. Gold CoS is USD
per attributable ounce sold. Copper prices and CoS are USD/lb, attributable
sales are kt, and $1~\mathrm{t}=2204.6226218488~\mathrm{lb}$. Values use an
explicit scale, normally USD millions; per-share values require a reconciled
diluted-share denominator. Already attributable volumes are never scaled by
ownership twice.

\subsection{Evaluation design}

Model fit, predictive performance, numerical convergence and economic sensitivity are separate questions. In-sample residual improvement is not evidence of out-of-sample superiority. Likewise, source-model parity is a prerequisite for integration, not a new economic finding.

Error metrics must use the scale relevant to the valuation. Team~4's
12-quarter cost holdout against a random walk reports RMSE $0.0504$ versus
$0.1512$ and MAPE $0.603\%$ versus $1.77\%$ on the log master chain
\cite{team4}. A log error of $0.05$ is roughly a 5\% level error; MAPE on
log values is not MAPE in USD/oz and depends on the log level and unit scale.
The RMSE ratio is the useful comparison, not the stated MAPE level.

Coverage and uncertainty sources must be explicit. The 95\% cost bands and
80\% production bands describe operating uncertainty; the frozen Chapter~12
quantiles vary gold prices and discount rates with production and costs fixed.
The copper extension adds copper-price uncertainty under declared dependence
scenarios. These quantities are neither interchangeable nor complete
corporate confidence intervals.

Copper reserved-strike scores test interpolation. Temporal scores fit only
origin-date quotes and project variance/intensity to the next admitted date;
target forward and as-of rates condition repricing. Common target support,
origin-mean and inside-hull persistence benchmarks, date-equal and pooled
losses are retained separately. Conditional IV repricing does not test a
physical price forecast or the accounting cash-flow bridge.

\subsection{Research boundaries}

This document is a reproducible working compilation. It contains no investment recommendation. Source figures and quantitative tables are withheld from the master until their asset, generator, data and recomputation status are accepted.
